% !Mode:: "TeX:UTF-8"

\chapter{引言}\label{ch:1}

\section{研究背景以及意义}

教育信息化推进之后，高校体育教学中的体质健康管理已从“纸质记录加人工汇总”逐步转向“系统录入加在线统计”。实际教学中，体测工作仍然面临数据量集中、录入周期短、项目规则固定、跨班级核对频繁等问题。教师和管理人员往往要在短时间内完成学生信息维护、体测成绩录入、结果核对、等级判定和后续分析，重复录入、漏录、口径不一致和追溯困难等情况依然常见。学生也需要查询个人历史成绩、课程安排和请假结果，传统分散式处理方式已难以满足当前高校体测管理的效率要求。

除效率问题外，体测业务还带有明显的规则约束。大学生体质评价需要遵循《国家学生体质健康标准》的评分依据\cite{national2014fitness}，系统不仅要“记得住数据”，还要“算得准成绩、查得到来源、说得清依据”。体测管理系统不能停留在简单的信息登记层面，应形成课程组织、数据录入、成绩计算、结果分析和日志追溯相互衔接的业务闭环。

基于上述背景，本文设计并实现了一个大学生体质健康智能分析系统。系统围绕课程、学生、体测记录、成绩结果和请假审批组织业务数据，将管理端、教师端和学生端纳入同一业务框架，为后续智能分析与结果解释提供基础支撑。

本文的主要创新点如下：

（1）提出基于业务数据约束的AI分析方法。系统在调用模型前先按学期、课程和学生范围整理统计数据或历史记录，将原始数据转换为结构化提示信息，模型输出受到真实业务数据约束，不脱离数据凭空生成结论。

（2）构建AI问答证据组织机制。AI问答不将原始数据直接送入模型，而是根据问题意图提取学生记录、课程统计、排行信息和条件样本等结构化证据，再与对话上下文组装为证据摘要，使回答可追溯、可复核。

（3）实现分析结果可追溯机制。AI分析与问答的输入摘要、输出文本和生成时间统一写入分析日志表，形成“业务数据整理—模型调用—结果留痕”的完整链路。

\section{国内外研究现状}

国外高校的体质健康管理系统起步较早，研究多集中在在线采集、健康评估、趋势跟踪和个性化干预等方向。部分高校已将学生体测数据纳入统一校园平台，通过标准化录入和自动统计减少人工汇总的工作量。一些研究还引入了运动表现分析、风险预测和健康建议生成等功能，用于观察学生体能变化并辅助教学决策。国外系统在平台建设和数据利用方面较为成熟，但其评价标准、课程组织方式和业务边界与我国高校体测制度并不一致，不能直接套用于国内场景。

国内高校体测管理系统近年来发展较快，许多系统已能完成学生档案维护、成绩录入、统计查询和报表导出等基础功能，部分平台也开始支持可视化分析和移动端访问。SpringBoot与Vue框架组合在校园管理系统中已有应用实践\cite{yao2025staffmanage}。但从实际使用效果看，仍存在几个问题。多角色权限边界不够清晰，教师、学生和管理员的操作范围容易交叉。批量导入、异常校验和课程状态控制不够细致，集中体测时数据一致性容易受影响。系统虽然能展示统计结果，但对结果形成过程和数据依据的解释能力有限，缺少面向真实业务数据的问答式分析能力。

为更清楚地说明国内外体测管理系统的差异和本系统的定位，表 \ref{tab:domestic-foreign-comparison} 从功能覆盖、权限控制、数据校验和智能分析等维度进行了对比。

\begin{table}[H]
  \centering
  \caption{国内外体测管理系统对比}
  \label{tab:domestic-foreign-comparison}
  \footnotesize
  \renewcommand{\arraystretch}{1.35}
  \begin{tabularx}{\textwidth}{|p{2.2cm}|p{3.0cm}|p{3.0cm}|p{3.0cm}|}
    \hline
    对比维度 & 国外典型系统 & 国内现有系统 & 本系统 \\
    \hline
    功能覆盖 & 在线采集、健康评估、趋势跟踪为主 & 学生档案、成绩录入、统计查询为主 & 课程管理、体测录入、成绩计算、请假审批、AI分析问答一体化 \\
    \hline
    权限控制 & 按角色分层，粒度较细 & 角色边界常不够清晰，存在越权风险 & 基于RBAC + 课程归属的三级权限隔离 \\
    \hline
    数据校验 & 标准化录入，校验较完善 & 批量导入校验不够细致 & 服务端多层校验：格式、范围、唯一性、课程状态 \\
    \hline
    智能分析 & 部分系统具备风险预测和个性化建议 & 多为统计图表展示，缺少解释能力 & AI分析基于业务数据约束，AI问答附带证据来源与日志追溯 \\
    \hline
    请假审批 & 通常独立于体测系统 & 与体测业务关联较弱 & 请假与课程绑定，特殊标签自动写回体测备注 \\
    \hline
  \end{tabularx}
\end{table}

针对上述不足，本文在常规体测管理功能之外，进一步引入学生端自助查询、请假审批与课程绑定、AI 辅助分析与问答、分析结果留痕和日志追溯等设计，使系统不仅能够完成数据管理，也能围绕“录入—计算—分析—解释—追踪”形成较完整的业务闭环。近年来，生成式人工智能技术在各领域的应用持续扩展\cite{guo2025chatgpt}，基于大语言模型的信息检索与分析方法也逐步应用于实际业务系统\cite{xie2024llm}，为体测管理系统引入智能分析能力提供了技术基础。

\section{论文研究内容}

\subsection{主要研究内容}

本文围绕高校体测管理业务完成系统设计与实现，重点解决课程、学生和成绩数据的组织方式，管理员、教师和学生三类角色的分工，以及 AI 能力如何嵌入既有业务流程等问题。系统以课程管理和体测数据处理为主线，覆盖登录认证、权限控制、批量导入导出、成绩计算、学生请假、日志记录、统计分析和智能问答等功能。

与传统体测管理系统相比，本文将课程绑定、请假审批、日志留痕和 AI 分析问答纳入统一流程，系统既能支持体测数据处理，也能围绕历史记录生成有依据的分析结论。

\subsection{采用的手段}

系统采用前后端分离架构。后端使用 SpringBoot 3.3 与 MyBatis-Plus 实现业务功能和数据访问，结合 SpringSecurity 构建 RBAC 权限模型。前端基于 Vue3 与 Element Plus 完成页面交互。文件处理方面，系统通过 Excel 模板支持批量导入导出，并在服务端完成格式、范围和唯一性校验。

\subsection{论文成果}

本文完成的主要成果包括：管理端支持用户、角色、课程、教师分配、成绩分析和日志审计等管理功能，能够完成基础数据和课程数据的批量导入，并在课程结束时自动进行成绩计算和状态控制。

教师端能够查看本人课程、录入或导入体测数据、处理学生请假并完成课程结果复核。学生端能够查询个人历史体测、课程安排并提交请假申请。系统同时加入 AI 分析与问答能力，用于辅助数据解读和结果核对。
